\subsection{The incoming rebased frame and two shorter centers}
The latest Nearest factorial input already changes the first-index
congruence and the magnitude budget together, obtaining degree 298 at
150 gates. Its rebased frame replaces the two old definitions by
\begin{equation}\label{eq:rebased-frame}
 k=R+1+h\ell,\qquad \mathcal B_6=Y^2B_0-k^2q^6,
 \qquad \ell=XY^2,
\end{equation}
and uses $D=\mathcal B_6-4$. Every other register in the preceding
complete definitions is retained. The older frame instead keeps
$k=R+1+hE$ and $\mathcal B_8=Y^2B_0-k^2q^8$.
These are different polynomial families. A degree-298 endpoint in the
rebased family does not contradict a degree-300 bound in the older one.

Set $\gamma=8\Delta+4$ and
$\lambda=4(\gamma U-rc^4)$. The new centers are
\begin{align}
 \Psi'_{298}&=\lambda V+U^2T
  \left\{rc^4[(\gamma-1)(c^2-2)-2]+64A\Delta\right\},
       \label{eq:center298}\\
 \Psi'_{300}&=\lambda V+U^2T
  \left\{rc^4[(\gamma-1)(c^2-2)]+64A\Delta\right\}.
       \label{eq:center300}
\end{align}
Use $\Psi'_{298}$ with the rebased front and $\Psi'_{300}$ with the
old front, in both cases with $L=32$ in \eqref{eq:scalarfinal}.
The ten outer coordinates and the eleventh coordinate $n$ remain
unrestricted integers. The six authentic numerical ports are unchanged.

\begin{theorem}\label{thm:center-shear}
The two displayed equations have polynomially bijective integer zero sets
with their respective calibrated parents. They have exact degree and live
binary-operation counts $298/147$ and $300/145$, respectively.
\end{theorem}
\begin{proof}[Proof of the zero set assertion]
The incoming rebased center is
\[
 \Psi^{\rm p}_{298}=[(32\Delta+16)U-4rc^4]V
 +U^2T[8U(c^2-2)+64A^3+rc^4(3c^2-16)].
\]
The old degree-300 center replaces its final term $rc^4(3c^2-16)$
by $3rc^6$. The exact identities
$U=rc^2+\Delta rc^4$ and $A^3=A\Delta+A$ give
\begin{align*}
 \Psi'_{298}-\Psi^{\rm p}_{298}
      &=U^2T(16rc^2-64A)=:\epsilon_{298},\\
 \Psi'_{300}-\Psi^{\rm p}_{300}
      &=U^2T(16rc^2-64A-14rc^4)=:\epsilon_{300}.
\end{align*}
Each $\epsilon$ is an integer polynomial independent of $n$.
Keeping all outer coordinates and setting $n_{\rm p}=n+\epsilon$
makes the parent's $Z$ exactly equal to the new $Z$, hence makes the
entire parent and new polynomials equal under substitution.
The inverse is $n=n_{\rm p}-\epsilon$. This is a bijection on the
full integer zero sets, not a rational map or a test at sample points.
It proves the new calibration once each respective parent is calibrated.
The operation and exact degree assertions are proved below.
\end{proof}

\subsection{The rebased parent keeps its guard and its late scale proof}
The changed frame in \eqref{eq:rebased-frame} must not be confused with
an earlier rejected guard deletion. The complete incoming ordinary proof
is retained with its dependencies. Its critical witness argument is stated
here to fix the logical handoff precisely.

The unchanged even-weight norm, coprime-product and adjacent-coefficient
descent recover $B=u^2$, $D_f=f^2$, $D_t=\tau^2$ and an integral
auxiliary $y$, with $u,f,\tau>0$, together with
$\mathcal B_6\ge12$, $S\ne0$ and $S\mid A_c$.
Before extraction one still has $\Delta>1$, odd nonzero $H$,
$D\equiv8$ or $12\pmod{16}$, $P\equiv1\pmod4$, and odd nonzero
$\delta$. The possible cancellation of two negative radical parts would
require $1024B=P$, excluded modulo four. The exact coprimalities
\[
 \gcd(B,D_f)=\gcd(B,\delta C_y)=\gcd(D_f,\delta C_y)=1,
 \qquad D_f\equiv1\pmod3
\]
therefore supply positive individual squares. With the proof-only values
$I=c^2u$ and $T_{\rm old}=uT$, the recovered equations are
\[
 f^2-\Delta I^2=1,\qquad
 V=f^2(c-R+4IT_{\rm old})-c,
\]
\[
 (\Delta I)^2V^2-[(\Delta I)^2-1]y^2=1.
\]
In particular $c^2\mid I$. The denominator descent uses this integral
quantity, not an assumed integral value of a rational square root.

The guard $\mathcal B_6>0$ gives $kY\ne0$, $B_0>0$ and
$|Y|>|q|^2$. Its centered disk supplies the common nonzero signs of
$q,F,C,q-F-C$, the strip $|c/k-Y|<1/2$, and the authentic packing bound
\[
 |q|(|q|^2-1)<R<|q|^4-2.
\]
Since $|X|\ge4|q|$ and $|q|\ge16$,
\[
 |\ell|=|X|Y^2>4|q|^5>3R+1.
\]
The carrier valuation gives $M\mid H(Y-q^3)$, whose 2-primary part
makes $Y$ even, and $H\mid q^{15}K_q^5$ for
$K_q=X(X+2ac)-1$. It does \emph{not} yet give $q^3\mid Y$.

The first Pell recurrence supplies
$\psi_{2\ell+1}(m)\equiv m\pmod{2\ell}$.
Together with $k=R+1+h\ell$ it gives the signed congruence
$\varepsilon\,2m\equiv R+1\pmod{|\ell|}$.
Its positive branch forces $m\ge(R+1)/2$; the negative branch forces
$m>R$. The retained absolute root folding and nonodd strong-rank lemma,
using $c^2\mid I$ rather than an unjustified oddness assumption, and
the two auxiliary residues give the main index $p=R$. The negative
branch is impossible, and no wrap in the positive branch gives
$2m=R+1$. The retained ratio argument needs $|Y|\ge6$, which has
already been proved; it does not need the as-yet unavailable cubic scale.
Writing $R=2t+1$, it yields
\[
 |Y|>(|X|-3)^t/2-1,
 \qquad |H|>|X|(|X|-3)^t.
\]

Put $x=|X|$ and
$\Pi=(\sgn(c)X-2^R)(2^R\sgn(c)X+1)$.
The inherited whole-product bound, applied with these proved hypotheses,
gives
\[
 H\mid q^{15}\Pi^5,\qquad
 |q|^{15}|\Pi|^5<2^{20t-10}x^{25},
 \qquad |H|>x^{t+1}/2^t.
\]
The authentic bounds $x\ge2^{64}$ and $t\ge36$ make
$43t-1526\ge22$, so the divisible integer has magnitude below $|H|$.
Thus $\Pi=0$ and $\sgn(c)X=2^R$. Reduction modulo three fixes
$c>0$ and $X=2^R$, and the surviving sign branch and strip give $k,Y>0$.
Only at this point does $q\mid X$ show that $|q|$ is a power of two.
Now odd $H$ is coprime to $M=2q^3$, and the scale is recovered late:
\[
 Y=(2v_{\rm scale}+1)q^3,\qquad v_{\rm scale}\in\Z.
\]
The centered half-binomial theorem then gives $0<\eta/k<1/4$.

A general rebased zero need not satisfy $\mathcal B_8>0$.
Soundness therefore returns to the earlier $B_0$-only source, not to a
stronger free-scale guard. In that native handoff set
$h_{\rm old}=Yh$. Then
$R+1+h_{\rm old}XY=R+1+hXY^2=k$ exactly.
The retained construction regenerates the required fresh odd strong and
auxiliary witnesses, restores its native interval and square coordinate,
and passes an actual source zero to its compiler/refutation theorem.
It does not assert that the current multiplier was already odd.

Conversely the canonical source has
$k=2\psi_{2\ell+1}((R+1)/2)$, whose recurrence implies
$4\ell\mid k-R-1$, so the new $h$ is integral.
Its $\mathcal B_8>64$ implies $\mathcal B_6>\mathcal B_8>64$.
Choose the Pell guard pair, then one strong pair satisfying
$4sAdc^2\mid I$, then $u,r$, then an intersective-quintic residue.
Interpolate $T_{\rm old}$ modulo $usr$ with that same fixed pair;
it is not reselected after imposing a residue. The retained divisibility
proofs make
\[
 n=A_c/S+32ufy+\tau g-\Psi
\]
an integer and make one norm factor vanish. This completes the identified
parent's witness construction. The exact center shear adds no further
number-theoretic assumption or witness-order choice.

\subsection{Exact degrees at the authentic coefficients}
For an independently checkable cancellation calculation, work first in the
free ring on $(A,c,r,T,R)$ with weights $(3,1,1,1,5)$, set
$\Delta=A^2-1$ and keep the literal $U,B,D_f,Q,V$. Define
\[
 D_\Psi=\Psi^2-1024BD_fV^2,
 \qquad E_\Psi=(Q-1)D_\Psi-1024BD_f(V^2-1).
\]
Complete exact collection gives the following full top forms:
\begin{align*}
 [D_{\Psi'_{298}}]_{80}
   &=512A^{15}c^{24}r^5TR-8192A^{17}c^{22}r^5T^2,\\
 [E_{\Psi'_{298}}]_{108}
   &=512A^{22}c^{30}r^6TR-24576A^{24}c^{28}r^6T^2,\\
 [D_{\Psi'_{300}}]_{82}&=512A^{16}c^{26}r^6T^2,\\
 [E_{\Psi'_{300}}]_{110}&=512A^{23}c^{32}r^7T^2.
\end{align*}
All higher-weight coefficients vanish. The supplied independent checker
reconstructs every coefficient, not only these displayed leading terms:
the two remainder pairs have $449/881$ and $455/888$ terms, respectively.

Put $W_0=Sn-A_c$. The exact cancellation identity is
\[
 p_a+p_b-p_c=\delta W_0^2+2\delta S\Psi W_0
              +\delta S^2P+S^2E_\Psi.
\]
Here $\deg S=20$, $\deg W_0=35$, $\deg\delta=28$ and
$\deg\Psi=54$. The first three terms have degrees at most
$98,137,122$ in the rebased frame, with the last bound $124$ in
the old frame. The final term has degree $148$ for $\Psi'_{298}$ and
$150$ for $\Psi'_{300}$. Thus the first inner square has degree $296$
in the rebased case, while $-4p_ap_b$ has exact degree $298$ and cannot
cancel. In the old frame the two contributions tie at degree $300$,
but the inner-square leader has $T$-degree four and the cross term has
$T$-degree two. Again they cannot cancel. Explicit actual coefficients are
\begin{align*}
 [j^{136}w^{46}v^{50}z^6s^6r^{12}c^{40}T^2]\,\mathcal N_{298}
   &=2^{272}3^{182}\beta^{136}>0,\\
 [j^{110}w^{50}v^{50}s^4r^{18}c^{64}T^4]\,\mathcal N_{300}
   &=2^{256}3^{160}\beta^{110}>0.
\end{align*}
Their exponent sums are $298$ and $300$. These are attained degrees at
the authentic fixed tuple, not generic-port degrees or sampled line tests.

\subsection{Two exact degree floors in two distinct fixed frames}
\begin{theorem}\label{thm:degree-floor}
Keep the eleven indeterminates, authentic tuple and all register definitions
of the respective scalar frame fixed. Among nonzero constant
integer radical weights $L$ and integer polynomial centers $\Psi$
independent of $n$, the minimum output degree is $298$ for the rebased
frame and $300$ for the old frame. The displayed centers attain those
bounds. Neither statement is an operation-count minimum or a bound for
other arithmetic frames.
\end{theorem}
\begin{proof}
Subscript zero now denotes the complete highest homogeneous part in the
actual coordinate ring. Put
\[
 q_0=6\beta j,\quad X_0=12q_0w,\quad a_0=X_0v,
 \quad F_0=12wC-q_0z,\quad R_0=-q_0^3F_0,
 \quad\ell_0=X_0v^2.
\]
For the rebased frame, $k_0=R_0+h\ell_0$; both terms have degree five.
For the old frame $k_0=R_0$. Set
\[
 q_*=ra_0^7c^6,\quad s_*=128sra_0q_0^{15},\quad
 p_*=s\ell_0^2k_0^3,
\]
\[
 G_0=(\mathrm{cone\_A}\,vF_0)^2+(2q_0v^2)^2+(q_0^3)^2.
\]
In the rebased frame,
$(\mathcal B_6)_0=-k_0^2G_0$ has degree $16$ and
$P_0=-p_*^2G_0$ has degree $54$.
In the old frame $P_0=-(q_0^4p_*)^2$ has degree $56$.
For both, $\delta_0=q_*$, $S_0=s_*$,
$\deg A_c=35$, and $\deg(BD_fC_y)=136$.
The complete forms are nonzero at the actual tuple. In particular
$G_0$ is a nonzero sum of three real polynomial squares since $q_0\ne0$.
The independent top-form checker includes the terms involving $h$ in
the rebased $k_0$ rather than incorrectly reusing the old $R_0$ alone.

Let $d=\deg\Psi$, with the zero center assigned $d=-\infty$.
If $d<54$, then $\deg Z\le73$ and $\deg p_a\le174$, while
$\deg p_c=176$ and $\deg p_b=122$ or $124$. Thus the inner sum
has degree $176$ and the cross term has degree at most $296$ or $298$,
so the output degree is $352$. If $d>54$, then
$\deg p_a=68+2d>176$ dominates, giving
$\deg\mathcal N=136+4d>352$; the cross term has smaller degree
$190+2d$ or $192+2d$.

For $d=54$, let $\psi_0$ be its nonzero degree-54 form. The rebased
cross term has complete degree-298 part
\[
 (-4p_ap_b)_{298}
       =(2q_*s_*^2\psi_0p_*)^2G_0.
\]
This is a nonzero sum of real polynomial squares. If the inner sum has
degree below $149$, this proves exact degree $298$; if above $149$,
its square has larger degree. At equality, the degree-298 part is this
nonzero sum plus the square of the inner sum's top form and cannot vanish.
The old frame has the identical argument at degree $300$, replacing
$G_0$ by $q_0^8$ and the dividing degree $149$ by $150$.
Thus both lower bounds hold. The attained coefficients above supply the
matching upper bounds.
\end{proof}
The theorem does not cover changed carriers, polynomial weights, other
radical architectures, altered outer-coordinate charts or arbitrary
Diophantine encodings. It neither assumes a zero nor decides consistency.

\subsection{Measured scalar results and the comparison boundary}
\begin{center}\small
\begin{tabular}{@{}lrrrr@{}}\toprule
Endpoint & Variables & Degree & Operations & Domain\\\midrule
Nearest factorial incoming &11&298&150&$\Z$\\
Round 19 control &11&300&148&$\Z$\\
New rebased center &11&298&147&$\Z$\\
New old-frame center &11&300&145&$\Z$\\\bottomrule
\end{tabular}
\end{center}
The degree-298 endpoint uses 29 additions, 24 subtractions and 94
multiplications; the degree-300 endpoint uses 29 additions, 23 subtractions
and 93 multiplications. Squares and fixed-scalar multiplications are
charged. The degree-298 endpoint improves both degree and cost over the
last delivered control, but its novelty against the latest input is a
three-gate saving, not a further degree reduction. The degree-300 endpoint
is a lower-cost tradeoff. Complete graph and coefficient payloads are
supplied; no expanded support or total coefficient height is asserted.

